\documentclass[11pt,a4paper]{article}
\usepackage[utf8]{inputenc}
\usepackage[french]{babel}
\usepackage[T1]{fontenc}
\usepackage{amsmath,amsfonts,amssymb}
\usepackage{geometry}
\geometry{margin=1.5cm}
\usepackage{tcolorbox}
\usepackage{multicol}

\title{\textbf{Fiche de Révision : Développement et Double Distributivité}}
\author{Mathématiques - Niveau 3\textsuperscript{e} (Série B)}
\maDate{}

\begin{document}

\maketitle

    \begin{tcolorbox}[colback=green!5!white,colframe=green!50!black,title=Méthode : La Double Distributivité]
    Pour développer $(a + b)(c + d)$, on effectue 4 multiplications :
    \begin{center}
        $(a + b)(c + d) = \mathbf{a \times c} + \mathbf{a \times d} + \mathbf{b \times c} + \mathbf{b \times d}$
    \end{center}
    \textbf{Attention aux signes :} Un nombre emporte toujours le signe qui est devant lui !
    \end{tcolorbox}
    
    \section*{EXERCICE 1 : Échauffement (Réduction)}
    \textit{Simplifier les expressions pour bien préparer le développement.}
    \begin{multicols}{2}
    \begin{enumerate}
        \item $A = 4x^2 - 5x + 2x^2 + 8x$
        \item $B = -3x + 7 - 2x - 10$
        \item $C = x^2 - 4x - 5x^2 + 4x$
        \item $D = -x + 5x^2 - 3 + 2x$
    \end{enumerate}
    \end{multicols}
    
    \section*{EXERCICE 2 : Double distributivité guidée}
    \textit{Compléter les pointillés pour développer les expressions.}
    
    \begin{enumerate}
        \item $E = (x + 3)(x + 5)$ \\
        $E = x \times \dots + x \times \dots + 3 \times \dots + 3 \times \dots$ \\
        $E = \dots\dots\dots\dots\dots\dots\dots\dots\dots\dots\dots\dots$
        
        \item $F = (x + 4)(x - 2)$ \\
        $F = x \times x + x \times (-2) + 4 \times x + 4 \times (-2)$ \\
        $F = \dots\dots\dots\dots\dots\dots\dots\dots\dots\dots\dots\dots$
    \end{enumerate}
    
    \section*{EXERCICE 3 : Maîtrise des signes}
    \textit{Développer puis réduire les expressions suivantes (attention aux nombres négatifs).}
    
    \begin{enumerate}
        \item $G = (x - 5)(x + 7) = \dotfill$
        \item $H = (x - 3)(x - 4) = \dotfill$
        \item $I = (2x + 1)(x + 6) = \dotfill$
    \end{enumerate}
    
    \newpage
    
    \section*{EXERCICE 4 : Vers le niveau Brevet}
    \textit{Développer et réduire ces expressions plus complexes.}
    
    \begin{enumerate}
        \item $J = (3x - 2)(2x + 5)$ \\
        $J = \dotfill$
        
        \item $K = (5x - 3)(x - 4)$ \\
        $K = \dotfill$
        
        \item $L = (-2x + 3)(x + 8)$ \\
        $L = \dotfill$
    \end{enumerate}
    
    \section*{EXERCICE 5 : Mélange distributivité simple et réduction}
    \textit{Développer d'abord la parenthèse, puis réduire avec le reste de l'expression.}
    
    \textit{Exemple : $2x + 3(x+1) = 2x + 3x + 3 = 5x + 3$}
    
    \begin{enumerate}
        \item $M = 5x + 2(3x - 4) = \dotfill$
        \item $N = x^2 + x(x - 5) = \dotfill$
        \item $O = 10 - (x + 2)(x + 1) = \dotfill$
    \end{enumerate}
    
    \section*{EXERCICE 6 : Géométrie et Algèbre}
    
    
    Soit un carré de côté $(x + 3)$.
    \begin{enumerate}
        \item Rappeler la formule de l'aire d'un carré : \dotfill
        \item Exprimer l'aire de ce carré en fonction de $x$ sous la forme d'un produit : \dotfill
        \item Développer et réduire cette expression pour obtenir l'aire sous forme "somme" : \\
        \dotfill
    \end{enumerate}
    
\end{document}